\section*{Capítulo \ref{ch:g9:genes-and-dna} --- Os genes e o DNA}

\begin{solution}{exo:g9:genes-and-dna:1}
Um \omterm{def:g9:genes-and-dna:gene}{gene} é um trecho de um \omterm{def:g9:chromosomes:chromosomes}{cromossomo} que carrega as
instruções de uma \omterm{def:g9:heredity-and-traits:traits}{característica}, num endereço fixo; os
\omterm{def:g9:genes-and-dna:gene}{alelos} são as versões dele --- o A, o B e o O do \omterm{def:g9:genes-and-dna:gene}{gene} do
grupo sanguíneo.
\end{solution}

\begin{solution}{exo:g9:genes-and-dna:2}
Porque os \omterm{def:g9:chromosomes:chromosomes}{cromossomos} vêm em pares: uma cópia de cada \omterm{def:g9:genes-and-dna:gene}{gene}
viaja em cada parceiro --- uma da mãe, outra do pai, no mesmo
endereço.
\end{solution}

\begin{solution}{exo:g9:genes-and-dna:3}
A ordem das letras químicas ao longo da fita. Quatro tipos de
letra, cerca de três bilhões delas num baralho humano.
\end{solution}

\begin{solution}{exo:g9:genes-and-dna:4}
As fitas são complementares letra por letra: separadas, cada
uma reconstrói a parceira, e uma escada vira duas escadas
idênticas --- uma cópia completa por célula-filha.
\end{solution}

\begin{solution}{exo:g9:genes-and-dna:5}
Um erro de cópia raro --- letras mudadas criando um \omterm{def:g9:genes-and-dna:gene}{alelo}
novo. Na maioria das vezes neutro, às vezes prejudicial, de
vez em quando útil.
\end{solution}

\begin{solution}{exo:g9:genes-and-dna:6}
Todos os \omterm{def:g6:exploring-our-environment:components}{seres vivos} escrevem e leem os \omterm{def:g9:genes-and-dna:gene}{genes} deles nas
mesmas letras e no mesmo código do \omterm{def:g9:genes-and-dna:dna}{DNA}. Consequências: a
prova mais forte de uma origem comum --- e o fato de
funcionar mudar \omterm{def:g9:genes-and-dna:gene}{genes} de uma \omterm{def:g5:biodiversity:species}{espécie} para outra, já que a
língua combina em toda parte.
\end{solution}

\begin{solution}{exo:g9:genes-and-dna:7}
Cada um dos pais carrega castanho com azul --- o castanho
expresso, o azul silencioso. Cada um distribui o azul com uma
chance em duas; as distribuições são independentes, então
azul com azul cai com uma chance em quatro --- o filho de
\omterm{def:g1:the-five-senses:sense}{olhos} azuis.
\end{solution}

\begin{solution}{exo:g9:genes-and-dna:8}
Grupo A: A com A, ou A com O. Grupo B: B com B, ou B com O.
Grupo AB: A com B --- os dois expressos. Grupo O: O com O.
\end{solution}

\begin{solution}{exo:g9:genes-and-dna:9}
Os dois têm a biblioteca inteira --- os mesmos 46 fios. O
\omterm{def:g8:nervous-communication:neuron}{neurônio} lê os \omterm{def:g9:genes-and-dna:gene}{genes} da maquinaria dos \omterm{def:g5:brain-in-command:system}{nervos} e deixa
fechados os capítulos da queratina da \omterm{def:g1:the-five-senses:sense}{pele}; a \omterm{def:g5:first-look-at-cells:cell}{célula} da \omterm{def:g1:the-five-senses:sense}{pele}
faz o contrário. A especialização é uma política de
empréstimo, e não uma biblioteca diferente.
\end{solution}

\begin{solution}{exo:g9:genes-and-dna:10}
Amasse a fruta (as \omterm{def:g5:first-look-at-cells:cell}{células} abertas); água salgada e
detergente (as \omterm{def:g6:cell-unit-of-life:parts}{membranas} --- os envelopes da \omterm{def:g5:first-look-at-cells:cell}{célula} e do
\omterm{def:g6:cell-unit-of-life:parts}{núcleo} --- dissolvidas, o \omterm{def:g9:genes-and-dna:dna}{DNA} solto e mantido \omterm{def:g7:oxygen-in-the-water:dissolved}{dissolvido});
coe (os restos fora); álcool gelado por cima (o \omterm{def:g9:genes-and-dna:dna}{DNA} sai da
solução no limite entre os dois): fios esbranquiçados que dá
para enrolar --- a própria molécula.
\end{solution}

\begin{solution}{exo:g9:genes-and-dna:11}
Nível da família: o queixo reaparece ao longo das linhas da
reprodução, carregado sem aparecer por uma geração. Nível dos
\omterm{def:g9:chromosomes:chromosomes}{cromossomos}: o determinante dele viaja num fio, reduzido à
metade nos \omterm{def:g8:sexual-reproduction:gametes}{gametas} e restaurado na \omterm{def:g5:flower-to-fruit:steps}{fecundação}, distribuído a
alguns filhos e a outros não. Nível da molécula: ele é um
\omterm{def:g9:genes-and-dna:gene}{alelo} --- uma grafia de algum \omterm{def:g9:genes-and-dna:gene}{gene} que dá forma ao rosto ---
copiado com fidelidade pelo século afora desde a \omterm{prop:g9:genes-and-dna:explains}{mutação} que
o escreveu pela primeira vez.
\end{solution}

\begin{solution}{exo:g9:genes-and-dna:12}
A \omterm{prop:g9:genes-and-dna:explains}{mutação} escreve: os erros de cópia criam grafias novas ---
nível da molécula, este capítulo. O embaralhamento
distribui: a redução à metade e a \omterm{def:g5:flower-to-fruit:steps}{fecundação} entregam a cada
filho uma combinação nova --- nível dos \omterm{def:g9:chromosomes:chromosomes}{cromossomos},
\cref{ch:g9:chromosomes} (e
\cref{prop:g8:sexual-reproduction:shuffle}). A expressão
mostra: das duas cópias distribuídas, o que é exibido e o que
viaja em silêncio --- a regra das duas cópias deste capítulo.
Só a \omterm{prop:g9:genes-and-dna:explains}{mutação} cria grafias genuinamente novas; as outras duas
recombinam e revelam o que a \omterm{prop:g9:genes-and-dna:explains}{mutação} escreveu um dia.
\end{solution}

\begin{solution}{pb:g9:genes-and-dna:1}
\textbf{1.} Uma \omterm{prop:g9:genes-and-dna:explains}{mutação}: um erro de cópia no \omterm{def:g9:genes-and-dna:gene}{gene} do
pigmento, uma vez, na linhagem de \omterm{def:g8:sexual-reproduction:gametes}{gametas} de algum
antepassado --- letras mudadas, as instruções estragadas, e a
grafia errada copiada com fidelidade desde então.

\textbf{2.} Porque as instruções funcionais de P constroem o
pigmento seja qual for o parceiro silencioso: p viaja
\emph{carregado sem aparecer} --- a expressão do álbum, agora
molecular.

\textbf{3.} Com p nos dois parceiros, não existem instruções
funcionais em nenhum dos dois endereços: a maquinaria do
pigmento --- o produto do \omterm{def:g9:genes-and-dna:gene}{gene} --- simplesmente nunca é
construída. A \omterm{def:g9:heredity-and-traits:traits}{característica} é a ausência de uma grafia
funcional.

\textbf{4.} Que os \omterm{def:g9:genes-and-dna:gene}{genes} de pigmento delas são
reconhecivelmente o mesmo \omterm{def:g9:genes-and-dna:gene}{gene} --- sequências parecidas no
mesmo código --- cópias aparentadas de um texto ancestral,
quebradas de forma independente na história de cada \omterm{def:g5:biodiversity:species}{espécie}.

\textbf{5.} As quatro distribuições: P com P, P com p, p com
P, p com p --- igualmente prováveis. A construção sem
pigmento: uma chance em quatro.

\textbf{6.} O pulo de geração: dois portadores que exibiam
pigmento transmitiram o que não mostravam, e as grafias
escondidas se encontraram --- o
\cref{ex:g9:heredity-and-traits:skip} rodado com o mecanismo
inteiro.

\textbf{7.} Todos carregam um p --- a contribuição inteira da
mãe deles naquele endereço é p. Nenhum mostra a condição: o P
do parceiro constrói o pigmento em todos eles.

\textbf{8.} O p da avó viajava num \omterm{def:g9:chromosomes:chromosomes}{cromossomo} de um par; a
redução à metade do \omterm{def:g8:sexual-reproduction:gametes}{gameta} dela distribuiu aquele parceiro
para o óvulo ou o \omterm{def:g8:reproductive-systems:male}{espermatozoide} que fez o pai ou a mãe
portador; a redução à metade desse pai ou dessa mãe o
distribuiu de novo para o \omterm{def:g8:sexual-reproduction:gametes}{gameta} que fez o neto --- fio,
\omterm{def:g8:sexual-reproduction:gametes}{gameta} e \omterm{prop:g8:sexual-reproduction:core}{zigoto}, duas vezes seguidas.

\textbf{9.} Família: dois pais de aparência comum, um filho
de uma chance em quatro, o carregar provado pelo \omterm{def:g2:born-grow-die:cycle}{nascimento}.
\omterm{def:g9:chromosomes:chromosomes}{Cromossomos}: os parceiros de um par carregando P e p,
reduzidos à metade e recombinados. Molécula: um \omterm{def:g9:genes-and-dna:gene}{gene}, um
\omterm{def:g9:genes-and-dna:gene}{alelo} funcional e um \omterm{def:g9:genes-and-dna:gene}{alelo} escrito errado, com a diferença de
grafia decidindo tudo.

\textbf{10.} De todo julgamento: em quem carrega, o defeito
de p é invisível --- o P do parceiro o cobre, então nada na
construção de quem carrega ``denuncia'' a grafia errada. O
que não dá para ver não dá para consertar e eliminar; as
grafias escondidas persistem indefinidamente.

\textbf{11.} A \omterm{def:g9:heredity-and-traits:traits}{característica} dela (nenhum pigmento) é
herdada --- p com p --- mas a \emph{consequência} dela é
vivida num \omterm{def:g6:exploring-our-environment:components}{ambiente} cheio de sol: o pigmento que falta era o
protetor solar do corpo, então o \omterm{def:g6:exploring-our-environment:components}{ambiente} escreve mais nela
do que nos pais protegidos --- a herança fixando a faixa, e o
mundo preenchendo-a.

\textbf{12.} O que copia: a fita dupla, com cada metade
reconstruindo a parceira a cada divisão celular e a cada
geração. O que foi copiado errado uma vez: algumas letras do
\omterm{def:g9:genes-and-dna:gene}{gene} do pigmento, num \omterm{def:g8:sexual-reproduction:gametes}{gameta} antigo. Por que o erro persiste:
a fotocopiadora é fiel --- ela copia o que estiver ali, e
copia as grafias erradas com a mesma fidelidade do sentido.
\end{solution}
